\begin{table}[htbp]
\centering
\caption{Complementary Outcomes: Which Margin Responds?}
\label{tab:jrap_complementary}
\small
\begin{tabular}{p{6.8cm}cc}
\toprule
Outcome & Coefficient (border effect) & $p$-value (cl.) \\
\midrule
log(revenue per active dispensary) & $+0.395$ & $0.015$ \\
log(revenue per capita) & $+0.046$ & $0.833$ \\
Number of active dispensaries (9-state entry-density panel) & null & n.s. \\
\midrule
\multicolumn{3}{l}{\textit{Placebo: general retail sales per capita, Colorado, 2016--2025}} \\
Border, full controls & $-0.006$ & $0.947$ \\
\bottomrule
\end{tabular}
\begin{minipage}{0.95\textwidth}
\vspace{0.3em}
\footnotesize
\textit{Notes:} The revenue-per-dispensary premium does not appear in total per-capita revenue or in the
number of active dispensaries, consistent with demand being absorbed by existing stores rather than by
entry, where entry is capped. The general-retail placebo uses Colorado Department of Revenue county-industry
filings for the 11 counties with disclosed industry detail \citep{CDORRetailByCounty}; the precise null
supports a cannabis-specific rather than a generic border-commerce interpretation.
\end{minipage}
\end{table}
